%% chapters/06-linalg-refresher.tex
\chapter{A Linear Algebra Refresher}
\label{ch:linalg}

\whereWeAre{%
We need eigenvalues and eigenvectors --- and we need them \emph{understood}, not just defined. This chapter rebuilds the linear algebra of symmetric matrices from scratch: eigenvalues, the spectral theorem, matrix norms, and the equivalence of all norms in finite dimensions.%
}

\PLACEHOLDER{Sources: existing main.tex spectral-theorem subsection; Gilbert Strang \emph{Introduction to Linear Algebra} for accessible-textbook tone.}

\section{Eigenvalues, eigenvectors, the spectral theorem}
\PLACEHOLDER{From scratch: $Av = \lambda v$, geometric meaning, characteristic polynomial; the spectral theorem $A = Q \Lambda Q^T$ for real symmetric $A$.}

\section{Quadratic forms and what eigenvalues control}
\PLACEHOLDER{The diagonal form $x^T A x = \sum \lambda_i y_i^2$; sign of eigenvalues controls definiteness.}

\section{Matrix norms}
\PLACEHOLDER{Frobenius and operator norms; submultiplicativity; equivalence of norms in finite dimensions.}

\section{Continuity of eigenvalues}
\PLACEHOLDER{Determinant is a polynomial in entries; characteristic polynomial coefficients are continuous; roots of polynomials depend continuously on coefficients; therefore eigenvalues depend continuously on matrix entries. This is the analytic engine of Chapter~9.}

\section{Brief warm-ups}
\PLACEHOLDER{3--5 short questions.}

\chapterRecap{%
\begin{itemize}
  \item \PLACEHOLDER{Spectral theorem: real symmetric $\Leftrightarrow$ rotation-scale-rotation.}
  \item \PLACEHOLDER{All norms equivalent in finite dimensions.}
  \item \PLACEHOLDER{Eigenvalues are continuous functions of the matrix.}
\end{itemize}%
}

\section*{Exercises}
\PLACEHOLDER{8--15 longer exercises.}
